\section{Scopo dell'esperienza}
usare lalgoritmo FFT ai dati sperimentali usando python (numpy.fft.rfft) per ottenere la magnitudine spettrale

\section{Dati da analizzare}
I segnali registrati dall'oscilloscopio digitale che mi sembrano siano quelli utili sono:
- Resistenza dinamica: Esperienza\_8/0.*\_CH1\_V.txt
- Raddrizzatore: Esperienza\_9/data.*\_CH*.txt
- Auto-oscillatore: Esperienza\_11/25.02.26.data*.txt
- Oscillatore smorzato RLC: Esperienza\_12/Pasqua.*.txt

\section{Struttura della relazione}
Per ogni set di dati si deve indicare i parametri sperimentali, mostrare il grafico FFT nel range di frequenze corretto e commentare brevemente i risultati

\section{Linee guida e Consegna, aka consigli del fuso per ricordarmi}
- Il documento finale deve essere un singolo PDF contenente i grafici e una chiara contestualizzazione
- inserire sempre i nomi degli assi, non usare il font matematico per le unità di misura e assicurarsi che la dimensione del testo sia "leggibile"
- Se svolto in gruppo, basta un solo upload una volta 

\section{Attività opzionali}
Filtro passa-banda per riduzione rumore o simulazione di lock-in numerico
